% !TEX program = xelatex
% =========================================================
% TEMPLATE DE MÉMOIRE LSM - UCLOUVAIN
% Adapté par Nicolas Nopenaire - Basé sur la structure EPL
% =========================================================
\documentclass{univ-thesis}

% --- Pour (dé)commenter ---
% 1. Sélection la(les) ligne(s) à (dé)commenter
% 2. Windows : Ctrl + Shift + /  & IOS : command + shift + /

% =========================================================
% 1. INFORMATIONS DU DOCUMENT
% =========================================================

\title{Titre principal du mémoire \\ (sur plusieurs lignes si nécessaire)}
\subtitle{Sous-titre éventuel (commenter la ligne si inutile)}

% Auteurs (Prénom en normal, NOM en petites capitales)
\author{Prénom \textsc{Nom}}
% \secondauthor{Prénom \textsc{Nom}}

% ---------------------------------------------------------
% Titre officiel du diplôme (Décommenter celui qui convient)
% ---------------------------------------------------------
\degreetitle{Master [120] en Ingénieur de gestion, à finalité spécialisée}
% \degreetitle{Master [120] en Sciences de gestion, à finalité spécialisée}
% \degreetitle{Master [60] en Sciences de gestion}
% \degreetitle{Master [120] in Business Engineering, professional focus}
% \degreetitle{Master [120] in Management, professional focus}

% ---------------------------------------------------------
% Majeure / Spécialisation (Exemples, voir README pour la liste complète)
% ---------------------------------------------------------
\myfocus{Majeure Business Analytics}
% \myfocus{Majeure Finance et Transition}
% \myfocus{Majeure Supply Chain Management}

% Promoteur / Superviseur
\supervisor{Prof. Prénom \textsc{Nom}}
% \secondsupervisor{Prof. Prénom \textsc{Nom}}

% Année académique
\years{2025--2026}


% =========================================================
% 2. PRÉAMBULE ADDITIONNEL (Optionnel)
% =========================================================
% % !TEX root = ../main.tex
% =========================================================
% PRÉAMBULE : PACKAGES ET CONFIGURATIONS AVANCÉES
% =========================================================

% --- Interligne ---
\usepackage{setspace}
\onehalfspacing % Interligne de 1.5 pour le confort de lecture

% --- Mathématiques ---
\usepackage{amsmath, amssymb} % Les standards pour les formules

% --- En-têtes et Pieds de page ---
\usepackage{fancyhdr}
\pagestyle{fancy}
\fancyhead{} % Nettoie l'en-tête par défaut
\fancyfoot{} % Nettoie le pied de page
\fancyfoot[C]{\thepage} % Numéro de page au centre

% --- Bibliographie (Norme APA) ---
\usepackage[backend=biber, style=apa, sorting=nyt, natbib=true]{biblatex}
% N'oublie pas de spécifier le chemin de ta bibliographie :
% \addbibresource{bib/references.bib}

% \usepackage{setspace}

\\begin{document}

% --- SÉLECTION DE LA LANGUE ---
% Le français est la langue par défaut.
% Pour rédiger en anglais, décommentez la ligne suivante :
% \selectlanguage{english}


% =========================================================
% 3. PAGES LIMINAIRES
% =========================================================
\maketitle

% % !TEX root = ../main.tex

% --- Résumé (Langue principale) ---
\chapter*{Résumé}
% Décommente la ligne suivante si tu veux que le résumé apparaisse dans la table des matières :
% \addcontentsline{toc}{chapter}{Résumé}

Ici, tu peux insérer le résumé de ton mémoire. Il doit synthétiser ta problématique, ta méthodologie, tes principaux résultats et tes conclusions en un ou deux paragraphes. C'est souvent la dernière chose que l'on rédige !

\vspace{1cm}
\noindent \textbf{Mots-clés :} Mot-clé 1, Mot-clé 2, Mot-clé 3, Mot-clé 4.

\vspace{2cm}

% --- Abstract (Version anglaise, optionnelle mais recommandée) ---
\begin{center}
    \Large \textbf{Abstract}
\end{center}
\vspace{0.5cm}

Here you can insert the English translation of your abstract. 

\vspace{1cm}
\noindent \textbf{Keywords:} Keyword 1, Keyword 2, Keyword 3, Keyword 4.
% % !TEX root = ../main.tex

\chapter*{Remerciements}
% \addcontentsline{toc}{chapter}{Remerciements}

La rédaction de ce mémoire marque la fin de mon parcours universitaire. À cette occasion, je tiens à exprimer ma profonde gratitude envers toutes les personnes qui ont contribué, de près ou de loin, à l'aboutissement de ce travail.

Je remercie tout d'abord mon promoteur, \textbf{Monsieur/Madame [Nom du promoteur]}, pour m'avoir guidé tout au long de cette recherche, pour ses conseils avisés et le temps qu'il/elle m'a consacré.

Je tiens également à remercier [Nom d'une autre personne, ex: un expert interviewé, un membre de la famille], pour [raison du remerciement].

Enfin, j'adresse toute ma reconnaissance à ma famille et à mes proches pour leur soutien indéfectible et leurs encouragements constants durant toutes mes années d'études.

\tableofcontents

% --- Test pour vérifier que la langue est bien celle désirée ---
% La langue actuellement comprise par LaTeX est : \languagename
\newpage

% =========================================================
% 4. CORPS DU DOCUMENT
% =========================================================
% \input{chapitres/01_introduction.tex}
% \input{chapitres/02_revue_littérature.tex}
% \input{chapitres/03_methodologie.tex}
% \input{chapitres/04_resultats.tex}
% \input{chapitres/05_conclusion.tex}

% =========================================================
% 5. BIBLIOGRAPHIE & ANNEXES
% =========================================================
% \bibliographystyle{apalike}
% \bibliography{bib/references.bib}

% Couverture arrière
\backcoverpage

\end{document}